\section{Experimental setup}
\label{sec:setup}

\subsection{Compute environment and its consequences}
All experiments were run on a single cloud virtual machine with four CPU cores of an Intel Xeon (2.1\,GHz, AVX-512 and AMX bfloat16 instructions), 15\,GB of memory and no GPU. PyTorch \citep{paszke2019pytorch} with bfloat16 autocast and channels-last tensors gave a $3$--$4\times$ speed-up over float32 on this hardware, which made the following design feasible: the segmentation network trains in about 23 minutes (30 epochs) and processes about eight $512\times512$ images per second at inference. Image embeddings come from a \emph{frozen} ImageNet-pretrained backbone, so that the grading stage trains in seconds and can be repeated over many seeds. This design imposes limitations that readers should keep in mind: (i) the segmentation network is trained at $512\times512$ rather than at native resolution (about $1950\times1950$), which shrinks microaneurysms to one to three pixels; (ii) it is trained on the 383 lesion-training images only; (iii) the image backbone is not fine-tuned on DR; (iv) segmentation results come from a single training run per configuration (seed 0), while grading-stage results average five seeds; and (v) heavy detectors (Faster R-CNN, DETR) and large transformer backbones could not be trained within the budget, so detection is evaluated through segmentation-derived proposals against a classical image-processing baseline and across segmentation losses.

\subsection{Implementation details}
\paragraph{Segmentation} Architecture, losses, sampling and optimisation are as described in Section~\ref{sec:seg}. Variants trained: (a) $\mathcal L_{\mathrm{bce}}$, (b) $\mathcal L_{\mathrm{bce+dice}}$, (c) $\mathcal L_{\mathrm{seg}}$ (focal Tversky) on raw colour input, and (d) $\mathcal L_{\mathrm{seg}}$ on Gaussian-filtered input (``GF''; Section~\ref{sec:prep}). All use the same seed, crop sampler and schedule. Thresholds for Dice/IoU are chosen on lesion-validation data per class.

\paragraph{Image embedding} ConvNeXt-T (ImageNet-1k weights \texttt{convnext\_tiny.fb\_in1k} from \texttt{timm}) is applied to $256\times256$ Gaussian-filtered images; global averages of stage 3 (384 channels) and stage 4 (768 channels) are concatenated to a 1152-dimensional vector and standardised using training-set statistics. We chose this backbone and resolution a priori for throughput reasons; no alternative was evaluated on the test set.

\paragraph{Grading heads} All heads are trained with AdamW (weight decay 0.01), cosine schedule over 60 epochs, mini-batches of 128, with learning rate $10^{-3}$ for MLPs and $5\times10^{-4}$ for the token head; the epoch with the highest quadratic-weighted kappa on validation half A (checked every three epochs) is retained. Validation is split into two halves with a fixed random permutation: half~A (1251 images) for model selection and temperature, half~B (1252 images) for conformal calibration. Each configuration is trained with five seeds (0--4); the data split is the same for all.

\subsection{Compared methods}
Table~\ref{tab:methods} lists the grading models. B1--B4 share the same embedding and (where applicable) the same lesion evidence as the proposed models, so that every difference is attributable to \emph{how} lesion information is used.

\begin{table}[!htbp]
\centering
\caption{Grading models compared. All use the same frozen ConvNeXt-T embedding $\varphi(x)$; ``evidence'' means the tensor $Z$ derived from the segmentation maps of the same segmentation network.}
\label{tab:methods}
\small
\begin{tabularx}{\textwidth}{@{}l>{\raggedright\arraybackslash}Xcccc@{}}
\toprule
Model & Description & Evidence & Ordinal & Soft & Rule \\
\midrule
B1 & MLP on $\varphi(x)$, 5-way cross-entropy & -- & -- & -- & -- \\
B2 & MLP on $\varphi(x)$, cumulative-link loss (Eq.~\ref{eq:ord}, $\sigma_y=0$) & -- & \checkmark & -- & -- \\
B3 & MLP on the flattened evidence $Z$ only & \checkmark & \checkmark & -- & -- \\
B4 & MLP on concatenation $[\varphi(x),Z]$ (naive fusion) & \checkmark & \checkmark & -- & -- \\
T0 & Evidence-token transformer, $\sigma_y=0$, $\lambda=0$ & \checkmark & \checkmark & -- & -- \\
T1 & T0 with soft targets ($\sigma_y=0.5$) & \checkmark & \checkmark & \checkmark & -- \\
T2 & T0 with rule loss ($\lambda=1$) & \checkmark & \checkmark & -- & \checkmark \\
\textbf{\mname} & transformer, $\sigma_y=0.5$, $\lambda=1$ & \checkmark & \checkmark & \checkmark & \checkmark \\
A1 & \mname\ without evidence tokens in the head (image token only; rule loss still reads $Z$) & (rule only) & \checkmark & \checkmark & \checkmark \\
\bottomrule
\end{tabularx}
\end{table}

For segmentation, the four U-Net variants of Section~\ref{sec:seg} are compared. For detection we compare segmentation-derived proposals from each variant and a classical baseline: top-hat filtering of the green channel followed by threshold and connected components (dark top-hat for MA and HE, bright top-hat for EX and SE), with the threshold tuned per class on lesion-validation data. We do not compare against published numbers of other systems because their splits, resolutions and training budgets differ.

\subsection{Evaluation metrics}
\paragraph{Grading} The primary metric is the quadratic-weighted Cohen's kappa (QWK) \citep{cohen1968kappa}, the standard agreement measure for ordinal DR grades. We also report accuracy, macro-averaged F1, the area under the ROC curve (AUC) for referable DR ($y\ge2$) computed from $P(y\ge2)$, expected calibration error (ECE, 15 bins), and the rule-violation rate (Section~\ref{sec:rule}). Conformal prediction is evaluated at $\alpha\in\{0.05,0.1,0.2\}$ by empirical coverage, mean set size, singleton fraction and referral sensitivity/specificity. All grading metrics are computed on the leak-controlled test set (Section~\ref{sec:data}) unless stated.

\paragraph{Segmentation} The area under the precision--recall curve (AUPR) of the pixel probabilities, computed exactly after quantising probabilities to 512 levels, is the primary metric because it is threshold-free and appropriate for extreme class imbalance \citep{saito2015pr}. We additionally report Dice \citep{dice1945} and intersection-over-union (IoU) at the per-class validation-optimal threshold, accumulated over all test pixels.

\paragraph{Detection} Average precision at IoU 0.5 and 0.3 (all-point interpolation \citep{padilla2021detmetrics}) and lesion-level FROC sensitivity at 1, 2, 4, 8 and 16 false positives per image, where a prediction is a hit if its centre lies in a ground-truth box. The relaxed IoU of 0.3 is reported because, for a box a few pixels wide, one pixel of displacement lowers IoU dramatically.

\subsection{Statistical analysis}
Grading results are reported as mean $\pm$ standard deviation over five seeds. For the main comparisons we additionally report 95\% confidence intervals obtained by a non-parametric bootstrap of the test images \citep{efron1993bootstrap} (2000 resamples) applied to the seed-averaged class probabilities, and paired bootstrap differences in QWK between models, which respect the pairing of test images. Differences are deemed significant when the 95\% interval of the paired difference excludes zero; no multiplicity correction is applied to the exploratory ablation table, and we say so where relevant. Segmentation metrics are reported from a single run per configuration; we therefore do not make significance claims among segmentation variants and we compare them descriptively.

\subsection{Reproducibility}
All scripts (data extraction, training, evaluation, figure and table generation) and intermediate artefacts (cached lesion maps, embeddings, per-seed probabilities) are released with the paper. Every number in the tables is generated by scripts from stored result files; there is no manual transcription. Random seeds are fixed. The code reads the DDR archive directly from its split zip volumes, so no re-packaging is needed.
